% !TeX root = ../../Book2.tex
% 纲要标题：### 12.3 不可分解模
% 纲要位置：工作记录/计划纲要.md:L1448

\section{不可分解模}
\label{b2:sec:1203}

合成列逐层取商，直和分解则要求每个分量在原模中有一个补模。两者的差别可由投影看清：一个直和分量对应一个满足 \(e^2=e\) 的自同态。本节先用自同态环描述不可分解性，再证明有限长度条件如何使模分解为自同态作用的幂零分量与可逆分量。

% 纲要要点（供目录核对）：
%
% * 局部自同态环的判别
% * Fitting 引理
% * 有限长度模上的幂等元分解
%

\subsection{局部自同态环的判别}
\label{b2:subsec:1203:01}

\begin{definition}\label{b2:def:indecomposable-module}
设 \(R\) 为含幺环，\(M\) 为非零幺性左 \(R\) 模。若不存在两个非零子模 \(A,B\) 使 \(M=A\oplus B\)，则称 \(M\) 为\term[bukefenjie-mo]{不可分解模}[indecomposable module]。即每次内部直和分解中，至少一个分量必须为零。
\end{definition}

简单模一定不可分解，但逆命题不成立。例如对素数 \(p\)，\(M=\mathbb Z/p^2\mathbb Z\) 的每个非零子群都包含唯一的阶 \(p\) 子群，所以两个非零子模不可能交于零。因此它不可分解，却不是简单模。这里并不否认其有两层合成列，而是说这两层不能在原模中分裂成直和。

\begin{corollary}\label{b2:prop:idempotent-direct-sum}
设 \(R\) 为含幺环，\(M\) 为非零幺性左 \(R\) 模。\(M\) 不可分解，当且仅当 \(\operatorname{End}_R(M)\) 中的幂等元只有 \(0,1\)。
\end{corollary}
\begin{proof}
\cref{b2:prop:idempotent-module-decomposition}已说明，幂等元 \(e\) 给出 \(M=\operatorname{im}e\oplus\ker e\)。若 \(e\ne0\)，像非零；若核为零，则 \(e\bigl(m-e(m)\bigr)=0\) 迫使 \(m=e(m)\) 对所有 \(m\) 成立，即 \(e=1\)。因此 \(e\ne0,1\) 恰好给出两个非零分量。反过来，两个非零分量的坐标投影显然既不为零，也不为恒等。
\end{proof}

比较两个直和分解时，会在某个分量的自同态环中得到恒等映射的有限和表达。若能从这个和中找出一个可逆项，就能比较它与另一分解中的某个分量。为此，我们需要一种环，使有限个非单位相加仍为非单位；这样它们的和便不可能是一。

\begin{definition}\label{b2:def:local-ring}
设 \(E\) 为非零含幺环。若 \(E\) 的全部非单位组成一个双边理想，则称 \(E\) 为\term[jubu-huan]{局部环}[local ring]。这里“单位”指具有两侧逆元的元素；不预设 \(E\) 交换。
\end{definition}

这里的局部性由环的非单位决定，与第六章的局部化构造是不同概念；局部化是一种使选定元素可逆的构造，并非每次局部化所得的环都是局部环。文献也常用“具有唯一极大左理想”定义局部环；在非交换情形，它与上述定义以及唯一极大右理想的条件等价。

\begin{proposition}\label{b2:prop:local-ring-criterion}
设 \(E\) 为非零含幺环。以下条件等价：\(E\) 为局部环；对每个 \(a\in E\)，\(a\) 与 \(1-a\) 中至少有一个为单位；\(E\) 有唯一极大左理想；\(E\) 有唯一极大右理想。在这些条件下，非单位组成的双边理想 \(J\) 是唯一极大左理想，也是唯一极大右理想，\(E/J\) 为除环。有限个非单位之和仍为非单位，特别地，若 \(1=x_1+\cdots+x_r\)，其中 \(r\ge1\)、\(x_i\in E\)，则至少一个 \(x_i\) 为单位。\(E\) 的幂等元只有 \(0,1\)。因此，若 \(R\) 为含幺环、\(M\) 为非零幺性左 \(R\) 模且 \(\operatorname{End}_R(M)\) 为局部环，则 \(M\) 不可分解。
\end{proposition}
\begin{proof}
若非单位构成理想 \(J\)，\(a\) 和 \(1-a\) 不能同时在 \(J\)，否则 \(1\in J\)，与一为单位矛盾。因此所述条件成立。

反过来，假设对每个 \(a\in E\)，\(a\) 或 \(1-a\) 为单位。先证明幂等元 \(e\) 只能为零或一：若 \(e\) 可逆，由 \(e^2=e\) 得 \(e=1\)；若 \(1-e\) 可逆，由 \(e(1-e)=0\) 得 \(e=0\)。这个限制还把一侧逆变成两侧逆。若 \(ab=1\)，则
\[
(ba)^2=b(ab)a=ba.
\]
幂等元 \(ba\) 不能为零，否则 \(b=bab=0\)，与 \(ab=1\) 矛盾。因此 \(ba=1\)，\(a,b\) 互为两侧逆。这一计算对只有左逆或只有右逆的情形都适用。

令 \(J\) 为非单位集合。若 \(a\in J\)、\(r\in E\)，\(ra\) 若可逆，则 \((ra)^{-1}r\) 是 \(a\) 的左逆；\(ar\) 若可逆，则 \(r(ar)^{-1}\) 是 \(a\) 的右逆。上段均会迫使 \(a\) 可逆，矛盾。所以 \(ra,ar\in J\)，并且 \(-a\in J\)。若 \(a,b\in J\) 而 \(u=a+b\) 可逆，则 \(u^{-1}a\) 与 \(1-u^{-1}a=u^{-1}b\) 都非单位，与假设矛盾。因此 \(a+b\in J\)，\(J\) 为不含一的双边理想。

对局部环的这个 \(J\)，任意非零剩余类由 \(a\notin J\) 代表，\(a\) 可逆，故 \(E/J\) 为除环。因此 \(J\) 在左右两侧均极大。任何真左理想或真右理想不能包含单位，故必包含于 \(J\)，所以它们各自的极大理想都只能是 \(J\)。

为从唯一极大左理想推出局部性，要用到每个真左理想都包含于一个极大左理想。对给定真左理想 \(L\)，考虑包含 \(L\) 的全部真左理想，按包含关系排序。它们组成非空偏序集；非空链的并仍是左理想，因为有限个元素可放进同一个链成员中检查减法与左乘，而这个并不含一，否则某个成员已经含一。空链以 \(L\) 为上界。由第一卷附录A定理A.1.4的 Zorn 引理，这个偏序集有极大元，即包含 \(L\) 的极大左理想。

现设 \(J\) 是 \(E\) 的唯一极大左理想。若 \(e^2=e\) 且 \(e\ne0,1\)，则 \(Ee\) 与 \(E(1-e)\) 都是真左理想：例如 \(re=1\) 会给出 \(e=(re)e=r e^2=1\)。由刚证的存在性，两者都包含于 \(J\)，便有 \(1=e+(1-e)\in J\)，矛盾。因此幂等元只有 \(0,1\)，前面关于 \(ba\) 的计算再次说明一侧逆必为两侧逆。

于是非单位 \(a\) 使 \(Ea\) 成为真左理想，必有 \(a\in J\)；反之 \(J\) 不含单位，所以 \(J\) 恰为非单位集合。它已对左乘封闭；若 \(a\in J\) 而 \(ar\) 可逆，则 \(a\) 有右逆，从而为单位，矛盾。因此 \(J\) 也对右乘封闭，是双边理想，\(E\) 为局部环。唯一极大右理想的情形在反环中应用同一论证。有限个非单位之和仍在理想 \(J\) 中，故不可能为单位；若它们之和为一，就至少有一项不在 \(J\) 中。幂等元结论与\cref{b2:prop:idempotent-direct-sum}还给出最后的不可分解性。
\end{proof}

例如，若 \(k\) 为域、\(n\ge1\)，则 \(k[t]/(t^n)\) 中元素可逆，当且仅当其常数项非零。确实，常数项非零的元素可写为 \(a(1-u)\)，其中 \(a\in k^\times\)、\(u\in(t)\) 幂零，逆由有限几何和给出；常数项为零的元素都在真理想 \((t)\)，不可能可逆。故非单位恰为 \((t)\)，这个环为局部环。但 \(\mathbb Z\) 不是局部环：非单位 \(2,-3\) 的和为单位 \(-1\)。虽然后者作为自身的模不可分解，反向局部性还需要下面的有限长度条件。

\subsection{Fitting 引理}
\label{b2:subsec:1203:02}

\begin{lemma}[Fitting]\label{b2:lem:fitting}
设 \(R\) 为含幺环，\(M\) 为有限长度幺性左 \(R\) 模，\(f\in\operatorname{End}_R(M)\)。对所有充分大的正整数 \(n\)，有
\[
M=\ker f^n\oplus\operatorname{im}f^n.
\]
在第一分量上 \(f\) 幂零，在第二分量上 \(f\) 可逆；两个分量都在 \(f\) 下稳定。可取任意 \(n\ge\max\bigl(1,\ell(M)\bigr)\)。
\end{lemma}
\begin{proof}
反复施加 \(f\) 产生两条链：
\[
\begin{aligned}
0&=\ker f^0\subseteq\ker f\subseteq\ker f^2\subseteq\cdots,\\
M&=\operatorname{im}f^0\supseteq\operatorname{im}f\supseteq\operatorname{im}f^2\supseteq\cdots.
\end{aligned}
\]
核记录逐次作用后变成零的元素，像记录仍能由反复作用得到的元素。有限长度同时限制严格升链与严格降链，所以两者都最终稳定。若某步 \(\ker f^r=\ker f^{r+1}\)，则对 \(f^{r+2}(x)=0\)，有 \(f(x)\in\ker f^{r+1}=\ker f^r\)，继而 \(x\in\ker f^{r+1}=\ker f^r\)。逐次归纳说明核链从此稳定。令 \(d=\ell(M)\)，从 \(\ker f^0=0\) 到 \(\ker f^d\) 若有一步相等，核链便从该步稳定；若前 \(d\) 步全都严格增加，则 \(\ell(\ker f^d)\ge d\)，所以 \(\ker f^d=M\)，此后也不再增加。因此从 \(r=d\) 起核链必稳定。由\cref{b2:prop:finite-length-endomorphism}，
\[
\ell(\operatorname{im}f^r)=\ell(M)-\ell(\ker f^r),
\]
故像链从同一位置起长度不变。相邻像有包含关系，长度相等时商长度为零，所以像也相等。

取满足上述界的 \(n\)，于是 \(\ker f^{2n}=\ker f^n\)、\(\operatorname{im}f^{2n}=\operatorname{im}f^n\)。要证明核与像交于零，考虑既经过 \(n\) 次作用、又会在 \(n\) 次作用后变成零的元素 \(y\)。写 \(y=f^n(x)\) 且 \(f^n(y)=0\)，就有 \(f^{2n}(x)=0\)。这正是比较 \(2n\) 与 \(n\) 的原因；由核稳定，\(x\in\ker f^{2n}=\ker f^n\)，从而 \(y=0\)。

对任意 \(x\in M\)，还要从像中减去一项，使余项落入核。由像稳定，\(f^n(x)\in\operatorname{im}f^{2n}\)，所以可取 \(z\in M\)，使 \(f^n(x)=f^{2n}(z)\)。减去 \(f^n(z)\) 后有
\[
f^n\bigl(x-f^n(z)\bigr)=0.
\]
因此 \(x-f^n(z)\in\ker f^n\)，而 \(f^n(z)\in\operatorname{im}f^n\)，得到所需分解。若 \(x\in\ker f^n\)，则 \(f^n(f(x))=f(f^n(x))=0\)，所以核在 \(f\) 下稳定；像也稳定，因为 \(f(\operatorname{im}f^n)=\operatorname{im}f^{n+1}=\operatorname{im}f^n\)。在核上 \(f^n=0\)；在像上，这个等式说明 \(f\) 满射，而 \(\ker f\subseteq\ker f^n\) 与已证的零交说明限制映射也单射，故可逆。
\end{proof}

这称为\term[Fitting-yinli]{Fitting 引理}[Fitting lemma]。\(\max(1,\ell(M))\) 是对所有自同态通用的指数上界，不一定是某个 \(f\) 的最小分解指数；\cref{b2:prop:fitting-integer-projections}中模的长度为五，乘二却在第三次作用时就得到分解。

\ProseParagraphBreak
对域 \(k\) 上有限维空间的算子 \(T\)，稳定核由那些被某个 \(T\) 的幂消去的向量组成，正是广义零特征子空间；稳定像是 \(T\) 可逆的部分。若最小多项式在 \(k\) 上分裂，\cref{b2:thm:jordan-canonical-form}的 Jordan 形把所有 \(J_r(0)\) 块放在第一部分，把非零特征值的块放在第二部分。Fitting 引理没有要求底环是域，也没有要求多项式分裂；有限长度提供的是核和像的同时稳定。

\begin{theorem}\label{b2:thm:indecomposable-local-end}
设 \(R\) 为含幺环，\(M\) 为非零幺性左 \(R\) 模。若 \(\operatorname{End}_R(M)\) 为局部环，则 \(M\) 不可分解。若 \(M\) 还有有限长度，则 \(M\) 不可分解，当且仅当 \(\operatorname{End}_R(M)\) 为局部环；在有限长度且不可分解情形，每个 \(f\in\operatorname{End}_R(M)\) 或者可逆，或者幂零，其全部非可逆自同态构成双边理想。
\end{theorem}
\begin{proof}
局部自同态环的幂等元只有零和一，由\cref{b2:prop:idempotent-direct-sum}，\(M\) 不可分解。这一方向没有使用长度假设。现设 \(M\) 有限长度且不可分解。对任意 \(f\)，Fitting 引理给出核与像的直和分解，不可分解性迫使
\[
\ker f^n=0\qquad\text{或}\qquad\operatorname{im}f^n=0.
\]
前一种情形使 \(f\) 单射，由\cref{b2:prop:finite-length-endomorphism}可逆；后一种情形使 \(f^n=0\)，即 \(f\) 幂零。

当 \(f\) 幂零时，\(1-f\) 可逆，因为若 \(f^n=0\)，有
\[
(1-f)(1+f+\cdots+f^{n-1})
=(1+f+\cdots+f^{n-1})(1-f)=1.
\]
因此每个 \(f\) 都使 \(f\) 或 \(1-f\) 可逆。\cref{b2:prop:local-ring-criterion}给出自同态环为局部环，同时保证非可逆元素对加法及左右乘法封闭。
\end{proof}

局部环一般只保证非单位组成理想，并不保证每个非单位都幂零。这里的幂零结论同时使用了有限长度与不可分解性；\cref{b2:prop:local-ring-nonnilpotent}将给出局部环中非幂零非单位的例子。

\ProseParagraphBreak

单有有限长度也不足以保证非可逆自同态幂零。例如对域 \(k\)，\(k^2\) 上的投影 \(\operatorname{diag}(1,0)\) 不可逆，也不幂零。相反，对循环 \(k[t]\) 模 \(k[t]/(p^n)\)，其中 \(p\) 不可约、\(n\ge1\)，任意两个非零子模均含有最小非零子模 \((p^{n-1})/(p^n)\)，所以它不可分解；其非可逆自同态因而都幂零。这与其自同态环 \(k[t]/(p^n)\) 中的非单位为 \((p)/(p^n)\) 相吻合。

\subsection{有限长度模上的幂等元分解}
\label{b2:subsec:1203:03}

\cref{b2:prop:idempotent-module-decomposition}也给出了有限版本：若自同态满足
\[
e_i^2=e_i,\qquad e_ie_j=0\ (i\ne j),\qquad
\sum_{i=1}^re_i=1,
\]
则 \(M=\bigoplus_i e_iM\)，反过来有限内部直和的坐标投影满足这些等式。现在进一步判别一个投影的像是否还能分解。

\begin{proposition}\label{b2:prop:orthogonal-idempotent-decomposition}
设 \(R\) 为含幺环，\(M\) 为幺性左 \(R\) 模。对非零幂等元 \(e\in\operatorname{End}_R(M)\)，其像 \(eM\) 不可分解，当且仅当 \(e\) 不能写成两个非零幂等元 \(u,v\) 之和且 \(uv=vu=0\)。若 \(M\) 有限长度，\(I\) 为集合，且 \((e_i)_{i\in I}\) 是 \(\operatorname{End}_R(M)\) 中满足 \(e_i\ne0\)、\(e_i^2=e_i\) 及 \(e_ie_j=0\ (i\ne j)\) 的一族元素，则 \(I\) 有限且 \(|I|\le\ell(M)\)。
\end{proposition}
\begin{proof}
若 \(e=u+v\) 如题述，则 \(eu=ue=u\)、\(ev=ve=v\)，所以两像都在 \(eM\) 中，并有 \(eM=uM\oplus vM\)，两项非零。反过来，若 \(eM=A\oplus B\) 且两项非零，先用 \(e\) 把 \(M\) 投到 \(eM\)，再投到 \(A,B\)，得到 \(u,v\)。它们分别在自身像上为恒等，在另一像及 \(\ker e\) 上为零，因而幂等、正交、非零，且 \(u+v=e\)。

对正交幂等元族，只需先看任意有限子集 \(F\subseteq I\)。若 \(x_i\in e_iM\) 且 \(\sum_{i\in F}x_i=0\)，施加 \(e_j\) 得 \(x_j=0\)，因为 \(e_j\) 在 \(e_jM\) 上为恒等，而在其余各项上为零。因此 \(\sum_{i\in F}e_iM\) 是内部直和。每个 \(e_i\ne0\) 使 \(e_iM\ne0\)，长度至少一，故由长度可加性，
\[
|F|\le\sum_{i\in F}\ell(e_iM)
=\ell\!\left(\bigoplus_{i\in F}e_iM\right)
\le\ell(M).
\]
若 \(I\) 有多于 \(\ell(M)\) 个元素，就能取到一个含 \(\ell(M)+1\) 个元素的有限子集，与此矛盾。所以 \(I\) 有限且满足所述界。整个论证只使用有限子族，没有对无限族定义幂等元之和。
\end{proof}

“正交”在这里指乘积为零，不涉及内积。若 \(M\) 有限长度，每个非零投影的像至少占一层长度，因此不能无限细分投影。如果 \(eM=A\oplus B\) 且两项非零，则 \(\ell(eM)=\ell(A)+\ell(B)\)，两个新分量的长度都严格小于原分量；继续分解时可对长度归纳。第十五章将把不能再写成两个非零正交幂等元之和的幂等元称为本原幂等元。

\ProseParagraphBreak
对内部直和 \(M=A\oplus B\)，自同态可写成由四个 Hom 群组成的分块矩阵；投影到 \(A\) 的幂等元是 \(\begin{pmatrix}1&0\\0&0\end{pmatrix}\)。在一般情形，它不在整个自同态环的中心，因为与含有非零交叉映射的分块矩阵不一定交换。模的直和分解需要的是幂等元，并不要求这些投影是中心幂等元；中心性对应更强的环论分解，后面另行讨论。

\begin{proposition}\label{b2:prop:fitting-integer-projections}
把 \(M=\mathbb Z/72\mathbb Z\) 看作幺性左 \(\mathbb Z\) 模，令 \(f(\bar x)=2\bar x\)。使 \(M=\ker f^n\oplus\operatorname{im}f^n\) 成立的最小正整数为 \(n=3\)，此时
\[
\ker f^3=\langle\bar9\rangle,\qquad
\operatorname{im}f^3=\langle\bar8\rangle.
\]
两个分量的投影分别为乘九与乘六十四，\(f\) 在第二分量上的逆为乘五。\(M\) 的长度为五。
\end{proposition}
\begin{proof}
\(f^3\) 为乘八，而 \(72=8\cdot9\)，故 \(8x\equiv0\pmod{72}\) 等价于 \(9\mid x\)。这给出所列核；像由 \(\bar8\) 生成。两分量分别有八与九个元素，交集为零，因为同时被八和九整除的整数被七十二整除。又 \(\bar9-\bar8=\bar1\)，所以它们共同生成 \(M\)，得到直和分解。这也就是中国剩余分解 \(\mathbb Z/72\mathbb Z\cong\mathbb Z/8\mathbb Z\oplus\mathbb Z/9\mathbb Z\)：乘二在前者上幂零，在后者上可逆。

乘九在 \(\langle\bar9\rangle\) 上为恒等，因为 \(9^2\equiv9\pmod{72}\)，在 \(\langle\bar8\rangle\) 上为零，因为 \(9\cdot8\equiv0\pmod{72}\)。因此它是第一分量的投影；另一分量的投影为 \(1-9=-8\equiv64\pmod{72}\)，即乘六十四。在阶九的第二分量上，乘二与乘五的复合为乘十，而 \(10\equiv1\pmod9\)，故乘五为所需逆。

当 \(n=1\) 时，\(\ker f=\langle\overline{36}\rangle\)，其中非零的 \(\overline{36}\) 也在 \(\operatorname{im}f=2M\) 中；当 \(n=2\) 时，\(\ker f^2=\langle\overline{18}\rangle\)，\(\overline{36}\) 又在 \(\operatorname{im}f^2=4M\) 中。因此前两步的核与像交集非零，最小指数确为三。最后，
\[
0\subsetneq\langle\overline{36}\rangle
\subsetneq\langle\overline{18}\rangle
\subsetneq\langle\bar9\rangle
\subsetneq\langle\bar3\rangle
\subsetneq M
\]
的相邻商依次为阶二、二、二、三、三的循环群，都是简单 \(\mathbb Z\) 模，所以这是长度五的合成列。
\end{proof}

\begin{proposition}\label{b2:prop:local-ring-nonnilpotent}
固定素数 \(p\)，令
\[
R=\mathbb Z_{(p)}
=\bigl\{\,a/b\in\mathbb Q\bigm|a,b\in\mathbb Z,\ p\nmid b\,\bigr\}.
\]
则 \(R\) 为交换局部环，非单位理想为 \(pR\)。其正则幺性左模 \({}_RR\) 不可分解，却不满足降链条件，因而不具有有限长度。非单位 \(p\in R\) 以及正则模上的乘 \(p\) 自同态都不幂零。
\end{proposition}
\begin{proof}
分母不被 \(p\) 整除的分数对加法、减法与乘法封闭，并含零与一，所以 \(R\) 是 \(\mathbb Q\) 的含幺子环，特别地为整环。对 \(a/b\in R\)，若 \(p\nmid a\)，其逆 \(b/a\) 仍在 \(R\) 中；若 \(p\mid a\)，则它在 \(R\) 中不可能有逆。否则存在 \(c/d\in R\) 满足 \(ac=bd\)，但左边被 \(p\) 整除、右边不被 \(p\) 整除，矛盾。因此非单位恰为 \(pR\)，这个集合是双边真理想，\(R\) 为局部环。由\cref{b2:prop:regular-hom-endomorphism}，\(\operatorname{End}_R({}_RR)\cong R^{\mathrm{op}}=R\)，故\cref{b2:prop:local-ring-criterion}推出正则模不可分解。

正则模的子模降链
\[
R\supsetneq pR\supsetneq p^2R\supsetneq\cdots
\]
处处严格：若 \(p^n\in p^{n+1}R\)，在分式域中消去 \(p^{n+1}\) 就得到 \(1/p\in R\)，矛盾。所以正则模不是 Artin 模，也不具有有限长度。由于 \(R\subseteq\mathbb Q\)，\(p^n\ne0\) 对每个 \(n\ge1\) 都成立；乘 \(p\) 自同态的第 \(n\) 次幂在元素一上取值为 \(p^n\)，同样非零。这给出局部自同态环中不幂零的非单位，而有限长度假设在这里恰好失败。
\end{proof}
